\scanpage{191}
donnant un paquet $\mathbb L$ de nbs premiers, et supposons que l'on a un \deleted{pro}groupe \inserted{profini} $\underline{\pi}_1^{\mathbb L}(X,e)$, relativement à une section $e$ de $X/S$ [plus généralement, un sous-schéma $\unsure{\underline e|\underline E}$ de $X$ fini, \deleted{et} radiciel surjectif sur $S$] Je suppose $X\mathrel{\text{---}}S$ plat de présentation finie, \inserted{à fibres géom\textsuperscript{t} connexes,} $S$ loc. noeth $(?)$ : la \inserted{cat. des} représentations linéaires de $\underline{\pi}_1^{\mathbb L}(X,e)$ dans des modules cohérents sur $S$ doit être équivalente à celle des \deleted{représentations} Modules coh. sur $X$ à strat. intégrables rel/$S$, qui sur chaque fibre \inserted{géométrique} correspondent à une représentation du groupe fondamental qui \deleted{soit} se fasse via un $\mathbb L$-groupe.

Plus généralement, si on a un \deleted{groupe \unsure{profini}} pro-(groupe \deleted{\unsure{étale|fini}} \inserted{étale}) \deleted{\unsure{loc. const.|loc. constant}} strict $\Pi$ sur $S$, et un pro-torseur $P$ sous $\Pi_X$ sur $X$, tel que \inserted{sur chaque} fibre géom \deleted{\unsure{$f_{\bar s}$|fibre}} $\Pi_{\bar s}$ apparaisse comme \uline{quotient} du groupe fondamental élargi de $X_{\bar s}$, alors on \deleted{doit} trouve une \margin{?} équivalence entre la catégorie des repr. (continues) de $\Pi$ \unsure{ds|dans} des modules coh. sur $S$, et celle des Modules coh. à stratifications intégrables sur $X$ qui, sur chaque fibre géom, veulent bien être subordonnées à $\Pi_{\bar s}$ (regardé comme quotient du groupe fondamental de $X_{\bar s}$,. \ldots
% archive sheet number: 190
